\section{Actores Responsables}

Los actores del sistema se dividen en tres categorías según su función dentro del ciclo de vida del mantenimiento:

\subsection{Roles Estratégicos y de Gestión}

\begin{itemize}
    \item \textbf{Gerente de Planta:} Usuario de alto nivel que consume información. Visualiza dashboards con Indicadores Clave de Rendimiento (KPIs) sobre la salud de los activos, costos y tiempos de inactividad para la toma de decisiones estratégicas \parencite{chen:2025, tractian:2025a, wireman:2013}.
    \item \textbf{Supervisor/Gestor de Mantenimiento:} Responsable de la estrategia general. Planifica los cronogramas de mantenimiento, gestiona presupuestos, aprueba, asigna órdenes de trabajo y supervisa al equipo operativo \parencite{accruent:2025, fracttal:2024, laurila:2017, tractian:2025b}.
    \item \textbf{Planificador de Mantenimiento:} Rol táctico que optimiza la programación de tareas, asegurando la disponibilidad de recursos como personal y repuestos para minimizar el impacto en la producción \parencite{fracttal:2024, laurila:2017, tractian:2025b}.
\end{itemize}

\subsection{Roles Operativos y de Campo}

\begin{itemize}
    \item \textbf{Técnico de Mantenimiento:} El usuario principal en campo. Utiliza una tablet para realizar inspecciones, ejecutar reparaciones, consultar documentación técnica y reportar el estado de las tareas.
    \item \textbf{Inspector HSEQ:} Especialista que utiliza la plataforma para realizar inspecciones de Activos Críticos de Seguridad (ACS), reportar anomalías y verificar su correcta resolución.
    \item \textbf{Jefe de Almacén:} Encargado de la gestión del inventario de repuestos y materiales de mantenimiento (conocido como inventario MRO). Procesa las solicitudes del sistema, actualiza el stock y registra la entrega de materiales.
    \item \textbf{Contratista Externo:} Usuario con acceso limitado que recibe órdenes de trabajo especializadas, ejecuta la labor y sube su informe para su procesamiento en el sistema \parencite{fracttal:2024, honeywell:2018, laurila:2017, servicenow:sf, tractian:2025b, wireman:2013}.
\end{itemize}

\subsection{Roles de Soporte Técnico y Analítico}

\begin{itemize}
    \item \textbf{Ingeniero de Proyectos (IT/OT/CAD):} Rol multifuncional responsable del Gemelo Digital. Se encarga de crear y mantener el modelo 3D (CAD), gestionar la infraestructura de software (IT) y asegurar la integración con sensores y maquinaria física (OT).
    \item \textbf{Analista de Datos y Fiabilidad:} Utiliza datos históricos y de sensores para desarrollar modelos predictivos, analizar tendencias de fallos, calcular métricas como el \textbf{Tiempo Medio Entre Fallos (MTBF)} y proponer mejoras en la estrategia de mantenimiento.
    \item \textbf{Administrador del Sistema:} Rol técnico que gestiona usuarios, permisos, integraciones con otros sistemas (ej. ERP) y realiza el mantenimiento general de la plataforma \parencite{chen:2025, ibm:2024, laurila:2017, mrzyk:2023, osullivan:2020, stum:2018, zhang:2025}.
\end{itemize}
